\documentclass[notes=show]{beamer}
%%%%%%%%%%%%%%%%%%%%%%%%%%%%%%%%%%%%%%%%%%%%%%%%%%%%%%%%%%%%%%%%%%%%%%%%%%%%%%%%%%%%%%%%%%%%%%%%%%%%%%%%%%%%%%%%%%%%%%%%%%%%%%%%%%%%%%%%%%%%%%%%%%%%%%%%%%%%%%%%%%%%%%%%%%%%%%%%%%%%%%%%%%%%%%%%%%%%%%%%%%%%%%%%%%%%%%%%%%%%%%%%%%%%%%%%%%%%%%%%%%%%%%%%%%%%
\AtBeginSubsection[
  {\frame<beamer>{\frametitle{Outline}   
    \tableofcontents[currentsection,currentsubsection]}}%
]%
{
  \frame<beamer>{ 
    \frametitle{Outline}   
    \tableofcontents[currentsection,currentsubsection]}
}
\usepackage{ulem}
\usepackage{pdflscape}
\setbeamertemplate{caption}{\insertcaption}
\usepackage{color,xcolor,ucs}
\usepackage{beamerthemesplit}
\usepackage{graphicx}
\usepackage{amsmath}
\usepackage{adjustbox}
\usepackage{pdfpages} 
\definecolor{mintgreen}{rgb}{0.6, 1.0, 0.6}
\makeatletter
\@addtoreset{subfigure}{framenumber}% subfigure counter resets every frame
\makeatother
\newcommand*{\Scale}[2][4]{\scalebox{#1}{$#2$}}%
\newcommand*{\Resize}[2]{\resizebox{#1}{!}{$#2$}}%
\usepackage{amsfonts}
\usepackage{booktabs}
\usepackage{color,soul}
\usepackage{pdfpages}
\setboolean{@twoside}{false}
\usepackage{amsmath}
\usepackage{pdflscape} 
\usepackage{eurosym}
\usepackage{amstext} % for \text
\DeclareRobustCommand{\officialeuro}{%
  \ifmmode\expandafter\text\fi
  {\fontencoding{U}\fontfamily{eurosym}\selectfont e}}
\usepackage[caption = false]{subfig}
\usepackage{appendixnumberbeamer} 
\usepackage{graphicx}
\usepackage{mathpazo}
\usepackage{hyperref}
\usepackage{multimedia}
\usepackage{graphicx}
\usepackage{multirow}
\usepackage{subfig}
\usepackage{verbatim}
\usepackage{booktabs, multicol, multirow}
\setcounter{MaxMatrixCols}{10}
\input{Macros.tex}
\AtBeginSection[]
{
  \begin{frame}[noframenumbering]
    \frametitle{Outline for section \thesection}
    \tableofcontents[currentsection]
  \end{frame}
}


\usepackage{xcolor,soul}
\renewcommand<>{\hl}[1]{\only#2{\beameroriginal{\hl}}{#1}}

% https://tex.stackexchange.com/questions/41683/why-is-it-that-coloring-in-soul-in-beamer-is-not-visible
\makeatletter
\newcommand\SoulColor{%
  \let\set@color\beamerorig@set@color
  \let\reset@color\beamerorig@reset@color}
\makeatother
\SoulColor
\AtBeginSection{\frame{\sectionpage}}

\newsavebox\mybox
\newenvironment{aquote}[1]
  {\savebox\mybox{#1}\begin{quote}\openautoquote\hspace*{-.7ex}}
  {\unskip\closeautoquote\vspace*{1mm}\signed{\usebox\mybox}\end{quote}}

\newtheorem{proposition}[theorem]{Proposition}
\newtheorem{Assumption}[theorem]{Assumption}
\newenvironment{stepenumerate}{\begin{enumerate}[<+->]}{\end{enumerate}}
\newenvironment{stepitemize}{\begin{itemize}[<+->]}{\end{itemize} }
\newenvironment{stepenumeratewithalert}{\begin{enumerate}[<+-| alert@+>]}{\end{enumerate}}
\newenvironment{stepitemizewithalert}{\begin{itemize}[<+-| alert@+>]}{\end{itemize} }
\usetheme{Madrid}
\makeatletter
\setbeamertemplate{footline}
{
  \leavevmode%
  \hbox{%
  \begin{beamercolorbox}[wd=.333333\paperwidth,ht=2.25ex,dp=1ex,center]{author in head/foot}%
    \usebeamerfont{author in head/foot}\insertshortauthor~~\beamer@ifempty{\insertshortinstitute}{}{(\insertshortinstitute)}
  \end{beamercolorbox}%
  \begin{beamercolorbox}[wd=.333333\paperwidth,ht=2.25ex,dp=1ex,center]{title in head/foot}%
    \usebeamerfont{title in head/foot}\insertshorttitle
  \end{beamercolorbox}%
  \begin{beamercolorbox}[wd=.333333\paperwidth,ht=2.25ex,dp=1ex,right]{date in head/foot}%
    \usebeamerfont{date in head/foot}\insertshortdate{}\hspace*{2em}
%    \insertframenumber{} / \inserttotalframenumber\hspace*{2ex} % DELETED
  \end{beamercolorbox}}%
  \vskip0pt%
}
\usepackage[style=british]{csquotes}

\def\signed #1{{\leavevmode\unskip\nobreak\hfil\penalty50\hskip1em
  \hbox{}\nobreak\hfill #1%
  \parfillskip=0pt \finalhyphendemerits=0 \endgraf}}

\makeatother
\usepackage{pdfpages}
\begin{document}
	\title[]{\large Minimum Wage		\vspace{0.5cm}
	}
	\author[Raffaele Saggio - Econ 560]{Raffaele Saggio}
	\institute[]{UBC}
	\date{\today \\}
	\maketitle

%%%%%%%%%%%%%%%%%%%%%%%%%	
%%%%%%%%%%%%%%%%%%%%%%%%%
%%%%%%%%%%%%%%%%%%%%%%%%%
%%%%%%%%%%%%%%%%%%%%%%%%%
\setbeamertemplate{footline}[frame number]{}

%gets rid of bottom navigation symbols
\setbeamertemplate{navigation symbols}{}

%gets rid of footer
%will override 'frame number' instruction above
%comment out to revert to previous/default definitions
\setbeamertemplate{footline}{}
\begin{frame}
\frametitle{Today}
\begin{itemize}
\item Minimum wage is at the heart of the labor demand theory
\bigskip
\item Plausibly exogenous change in price $\rightarrow$ what happens to quantity demanded.
\bigskip
\item Dogma of neo-classical theory: Demand curve slopes down!
\bigskip
\pause
\item Does it really?
\bigskip
\pause
\item Minimum wage literature instrumental for (i) ``Credibility Revolution" (ii) moving away from the perfectly competitive assumption of labor markets.  
\end{itemize}
\end{frame}
%%%
%%
\begin{frame}
\frametitle{Introduction}
\begin{itemize}
\item How do employers in low-wage labor markets respond to an increase in minimum wage?
\medskip
\item ``Conventional" Economic Theory is unambiguous: cut employment (Stigler, 1946).
\medskip
\item Time series regressions (teenager employment on minimum wages).
\medskip
\item Aggregate data seems to confirm this prediction (Brown 1982, 1983; Wellington, 1991). 
\end{itemize}
\end{frame}
%%
\begin{frame}
\frametitle{Wellington (1991)}
\begin{figure}[htb]
\centering
\includegraphics[width=1.1\textwidth]{aux/well}
\end{figure}
\end{frame}
\begin{frame}
\frametitle{The Right Minimum Wage}
\begin{figure}[htb]
\centering
\includegraphics[width=0.8\textwidth]{aux/right_mw}
\end{figure}
\end{frame}

%%
\begin{frame}
\frametitle{The power of simplicity and transparency}
\begin{figure}[htb]
\centering
\includegraphics[width=0.9\textwidth]{aux/quote}
\end{figure}
\end{frame}
%%%
\begin{frame}
\frametitle{Card and Krueger (1994)}
\begin{itemize}
\item 410 fast-foods restaurants in New Jersey and Pennsylvania following the increase in NJ minimum wage from 4.25\$ to 5.05\$ per hours.
\bigskip
\item Two research designs:
\bigskip
\begin{itemize}
\item Difference in Differences (compare NJ and PA before and after law change).
\medskip
\item Gap Design. (compare initially low wage establishments to high wage establishments before and after the design).
\end{itemize}
\bigskip
\item The value of surveys and ``zooming" in $\rightarrow$ employment, attrition, etc. 
\end{itemize}
\end{frame}
%%
\begin{frame}
\frametitle{Two Waves}
\begin{figure}[htb]
\centering
\includegraphics[width=0.9\textwidth]{aux/ck0}
\end{figure}
\end{frame}
%%
\begin{frame}
\frametitle{Always make sure to measure your treatment reliably!}
\begin{figure}[htb]
\centering
\includegraphics[width=0.9\textwidth]{aux/ck1}
\end{figure}
\end{frame}
%%
\begin{frame}
\frametitle{Ladies and Gentlemen we have a first stage!}
\begin{figure}[htb]
\centering
\includegraphics[width=0.9\textwidth]{aux/ck2}
\end{figure}
\end{frame}
%%
\begin{frame}
\frametitle{PA fast food restaurants is a good control group?}
\begin{figure}[htb]
\centering
\includegraphics[width=0.9\textwidth]{aux/ck3}
\end{figure}
\end{frame}
%%
\begin{frame}
\frametitle{Things look quite different in second wave (after law is passed)}
\begin{figure}[htb]
\centering
\includegraphics[width=0.9\textwidth]{aux/ck4}
\end{figure}
\end{frame}
%%
%%
\begin{frame}
\frametitle{Research Design \#1}
\begin{itemize}
\item Gap Design
\be
\Delta E_{i} = a + b X_{i} + c Gap_{i} + e_{i}
\ee
where 
\be
Gap_{i} = NJ_{i}\times \max\left(0,\dfrac{5.05-W_{1i}}{W_{1i}}\right)
\ee
and $W_{1i}$ is the reported starting wage of store $i$.
\bigskip
\item Questions for you:
\begin{enumerate}
\bigskip
\item What is $c$ measuring? (Mean value of $Gap_{i}=0.10$). 
\bigskip
\item Do I need PA fast-food restaurants to run this regression?
\bigskip
\item Is this a market level response?
\end{enumerate}
\end{itemize}
\end{frame}
%%
\begin{frame}
\frametitle{Research Design \#2}
\begin{itemize}
\item Difference in Differences
\be
E_{it} = \lambda_{t} + \gamma NJ_{i} + \delta ( NJ_{i}\times Post_{t}) + r_{it}
\ee
\item Card and Krueger estimate the model above in first differences (and adds controls).
\begin{enumerate}
\bigskip
\item Can you pin down $\delta$ from the previous table? 
\bigskip
\item What is the key identifying assumption here? What can go wrong?
\bigskip
\item Is $\delta$ identifying a market level response?
\end{enumerate}
\end{itemize}
\end{frame}
%%
\begin{frame}
\frametitle{Parallel Trends}
\begin{figure}[htb]
\centering
\includegraphics[width=0.9\textwidth]{aux/parallel}
\end{figure}
\end{frame}

%%
%%
\begin{frame}
\frametitle{Parallel Trends}
\begin{figure}[htb]
\centering
\includegraphics[width=0.9\textwidth]{aux/parallel2}
\end{figure}
\end{frame}
%%
\begin{frame}
\frametitle{Zero or positive?}
\begin{figure}[htb]
\centering
\includegraphics[width=0.9\textwidth]{aux/ck5}
\end{figure}
\end{frame}
\begin{frame}
\frametitle{More expensive Big Macs?}
\begin{figure}[htb]
\centering
\includegraphics[width=0.9\textwidth]{aux/ck6}
\end{figure}
\end{frame}
\begin{frame}
\frametitle{Modern Version of This}
\begin{figure}[htb]
\centering
\includegraphics[width=0.9\textwidth]{aux/orley}
\end{figure}
\end{frame}
\begin{frame}[plain]{The power of zero}

\begin{itemize}
\item A carefully thought out and transparent demonstration of the fragility
of min wage results

\begin{itemize}
\item Inferential issue: no clustering \textendash{} but shouldn't affect
GAP design
\item Results a bit under-powered to detect clear positive but strong enough
to reject big negative
\end{itemize}
\item Initial reaction of labor economists was not welcoming

\begin{itemize}
\item Hamermesh (1995): employers anticipated the change in ``before''
period and ``after'' period too soon for adjustment to occur.
\item Neumark and Wascher (2000): results driven by measurement problems
(responded to by Card and Krueger, 2000)
\end{itemize}
\end{itemize}
\end{frame}

\begin{frame}[plain]{The Backlash}

\begin{itemize}
\item {\footnotesize{}Businessweek: ``A Minimum Wage Study with Minimum
Credibility'' }{\footnotesize\par}
\end{itemize}
\begin{quotation}
{\footnotesize{}``Political correctness seems to have crept into
the inner sanctum of the AEA, discrediting its scholarly journal and
debasing its top prize. Unless the association cleans up its act,
it can kiss its credibility good­bye''}{\footnotesize\par}
\end{quotation}
\begin{itemize}
\item {\footnotesize{}James Buchanan in the WSJ}{\footnotesize\par}
\end{itemize}
\begin{quotation}
{\footnotesize{}``Just as no physicist would claim that 'water runs
uphill,' no self-respecting economist would claim that increases in
the minimum wage increase employment. Such a claim, if seriously advanced,
becomes equivalent to a denial that there is even minimal scientific
content in economics, and that, in consequence, economists can do
nothing but write as advocates for ideological interests.''}{\footnotesize\par}
\end{quotation}
\begin{itemize}
\item {\footnotesize{}Merton Miller in the WSJ}{\footnotesize\par}
\end{itemize}
\begin{quotation}
{\footnotesize{}``Raising the minimum wage by law above its market
determined equilibrium, they argue, actually costs nobody anything.
(Or at worst, costs nobody very much because it's only a small, marginal
increment, after all.) Is all this too good to be true? Damn right.
But it sure plays well in the opinion polls. I tremble for my profession.''}{\footnotesize\par}
\end{quotation}
\end{frame}

\begin{frame}[plain]{Interpretation}

\begin{itemize}
\item What to make of these results?

\begin{itemize}
\item Card-Krueger argue that positive employment effects reflect that not all workers are paid ``the margin" but firms have latitude to set wages (like in monospony models).
\end{itemize}
\item This is the key punchline of CK! Took many years to labor economists to digest this.
\end{itemize}
\bigskip{}

\begin{itemize}
\item Do GAP design and diff in diff identify the same parameter?
\begin{itemize}
\item Diff in diff measures market-wide response
\item GAP measures effect of raising wage on a single firm holding market
constant
\end{itemize}
\end{itemize}
\end{frame}

\begin{frame}[plain]{Harasztosi and Lindner (2018)}

\begin{itemize}
\item US min wage variation tends to be small and short run in nature
\item Hungary experienced a large (60\%) and persistent (\textasciitilde 8
years) increase in min wage in 2001
\end{itemize}
\end{frame}

\begin{frame}[plain]{Now we're talking..}

\includegraphics[scale=0.4]{aux/attila1}
\end{frame}
\begin{frame}[plain]{Harasztosi and Lindner (2018)}
\begin{itemize}
\item One reform, one country: need a control group!
\medskip
\item Lesson learned from CK: use Gap design!
\pause
\item Admin data from Hungary $\rightarrow$ can do Gap design on ``steroids"  
\item Use firm level exposure design to infer MW effects
\be
\dfrac{y_{it}-y_{i2000}}{y_{i2000}} = \lambda_{t} + \beta_{t} FA_{i} + r_{it}
\ee
where $FA_{i}$ is fraction of workers affected in firm $i$, i.e. the share of workers in firm $i$ being paid below the new minimum wage.
\medskip
\item What are the parametric assumptions of this model? 
\medskip
\item What would do with dying firms?
\medskip
\item What would do with new firms?
\end{itemize}
\end{frame}


\begin{frame}[plain]{Firm exposure  lowers employment}
\begin{center}
\includegraphics[scale=0.6]{aux/attila2}
\par\end{center}
\end{frame}
\begin{frame}[plain]{Is this effect large or small?}
\begin{center}
\includegraphics[scale=0.6]{aux/attila2}
\par\end{center}
\end{frame}
\begin{frame}[plain]{Firm exposure increases wages and total labor costs}
\begin{center}
\includegraphics[scale=0.6]{aux/attila3}
\par\end{center}
\end{frame}
\begin{frame}[plain]{Can you back up an elasticity of labor demand from these graphs?}
\pause
\begin{center}
\includegraphics[scale=0.3]{aux/attila4}
\par\end{center}
\end{frame}
\begin{frame}[plain]{How do employers respond?}
\pause
\begin{center}
\includegraphics[scale=0.4]{aux/attila5}
\par\end{center}
\end{frame}
\begin{frame}[plain]{Higher Q or Higher P? Use data on Price indexes for Manufacturing firms}
\pause
\begin{center}
\includegraphics[scale=0.4]{aux/reveneu_price}
\par\end{center}
How do you think this graph would look for other sectors (i.e. services)?
\end{frame}
%%
\begin{frame}[plain]{Employers in tradable sector responds differently...}
\begin{center}
\includegraphics[scale=0.4]{aux/attila8}
\par\end{center}
\end{frame}
\begin{frame}[plain]{Employers in tradable sector responds differently...}
\begin{center}
\includegraphics[scale=0.4]{aux/attila9}
\par\end{center}
\end{frame}
\begin{frame}[plain]{Who pays for the minimum wage then?}
\pause
\centering
\includegraphics[scale=0.5]{aux/attila_deco}
\end{frame}
\begin{frame}[plain]{Who pays for the minimum wage then?}
\begin{center}
\includegraphics[scale=0.5]{aux/attila7}
\par\end{center}
\end{frame}


\begin{frame}[plain]{Thoughts}

\begin{itemize}
\item Negative but small employment elasticity $\rightarrow$ out of 300,000 min wage workers in Hungary 30K lost their job while 270K had their wage rise by 60\%. 
\medskip
\item Cost effects of min wage largely passed through to consumers!
\item In an economy where min wage workers produced goods consumed by ``rich" people, then MW effectively represents a redistribution from rich to poor. 
\pause
\bigskip
\item Attila uses a neoclassical labod demand model (but with imperfect competition on the output market) to ``fit" its reduced form effects.  
\bigskip
\begin{itemize}
\item Unclear how to distinguish effects of \emph{firm}-specific from aggregate
price shocks at the market level. 
\item Monopsony model says small and big MW changes can have different signed
effects
\end{itemize}
\bigskip
\pause{}
\item Policy implications

\begin{itemize}
\item Price increases to domestic consumers are somewhat  problematic
\item Who is consuming these goods?
\end{itemize}
\end{itemize}
\end{frame}
\begin{frame}
\frametitle{Who pays the minimum wage?}
\begin{center}
\includegraphics[scale=0.5]{aux/trump_MC.jpg}
\par\end{center}
\end{frame}
\begin{frame}[plain]
\begin{center}
\includegraphics[scale=0.3]{aux/review_MW}
\par\end{center}
\end{frame}
\begin{frame}
\frametitle{More recent evidence from US Tax Returns of Independent Businesses: Wage Own Elasticity of around -0.25}
\begin{center}
\includegraphics[scale=0.2]{aux/max}
\par\end{center}
\end{frame}


%%
\begin{frame}[plain]{Dustmann, Lindner, Schonberg, Umkehrer and vom Berge (2021)}
\begin{itemize}
\item National wide introduction of the MW in Germany in 2015: 8.50 per hour (affecting around 15\% of workers). 
\medskip
\item High-quality matched employer-employee data 
\medskip
\item Gap design at individual + region level (they have ``high-quality" panel data) 
\medskip
\item Null effect on employment. But large increase in wages for affected workers.  
\medskip
\item Key new fact: ``reallocation effect" of the minimum wage. 
\end{itemize}
\end{frame}
%%
\begin{frame}[plain]{17\% of the wage increase following MW due to Reallocation}
\begin{center}
\includegraphics[scale=0.17]{aux/dustmann}
\par\end{center}
\end{frame}
%%
%%
\begin{frame}[plain]{What explains these positive re-allocation effects?}
\begin{itemize}
\item Original motivation to introduce MW in New Zealand (1890): ``helping
worthwhile companies that employ working class breadwinners against sweatshops that gained market shares by undercutting those worthwhile companies"
\medskip
\item Same idea for Scandinavian model: Collective bargaining will drive low performing establishment out of the market. 
\medskip
\item Underlying idea: firms have market power to set wages $\rightarrow$ wages depend upon productivity of the employer $\rightarrow$ MW kicks out ``bad" firms. 

\end{itemize}
\end{frame}
\end{document}

